% =====================================================================
% Table 2 (FIRST DRAFT) --- the two primitive questions and their answers
% ---------------------------------------------------------------------
% For Section 3 ("Two primitive questions"). The rows are the section's two
% questions; the columns the two poles of each.
% Alternative shape if this is not what you meant: a table of the *classic
% concepts* named in Section 2 para. 1 (Keynesian cross / Haavelmo /
% crowding out) mapped onto the closure rows.
% =====================================================================
\begin{table}[htbp]
\centering
\small
\begin{tabular}{@{}p{3.0cm} p{4.4cm} p{4.4cm}@{}}
\toprule
 & \textbf{Scarcity / clearing} & \textbf{Quantity adjustment} \\
\midrule
\textbf{Vision of the short run}
  & factors fully employed; prices and wages clear; output is
    supply-determined
  & output responds to demand; some price or quantity fails to clear \\[2pt]
\textbf{Theory of money}
  & immediate refinancing: every increase is offset elsewhere
  & no immediate budget constraint: external or endogenous financing \\[2pt]
\textbf{Classic reference}
  & crowding out; the balanced-budget (Haavelmo) theorem
  & the Keynesian cross; the multiplier \\
\bottomrule
\end{tabular}
\caption{The two primitive questions and their polar answers.}
\label{tab:two-questions}
\end{table}
